% IIPEM 2026, Section V (Results). Draft 2, 2 October 2026 (cut to the 2.5-page budget from draft 1).
% Source: results/iipem/ from analyze_iipem.py 6508969a6523ccc0 (figures only changed from 9fd80726350ee812;
%   every CSV and summary.md identical). summary.md 0bcf61787b416493 (sed 3d: df074d7173d18878),
%   paired_tests.csv 9fad719ba93bc2af, summary_iqm.csv c0d99fb395f8331a, within_arm_tests.csv 25f7ba87eadcce42.
% Tables I and II are generated by make_tables.py from the CSVs; do not edit their numbers by hand.
% The statistics and measurement definitions moved to methods_statistics.tex (\label{sec:stats}).
% \ref{sec:conditions} must point at the section that states the C1/C2 predictions.

\section{Results}\label{sec:results}

\begin{figure*}[!t]
\centering
\includegraphics[width=\textwidth]{fig_timeseries}
\caption{Per-step means over the 50 episodes (10 seeds $\times$ 5) of each arm; the vertical line is the event at step 50. Left: probe success $S$; the drop at step 50 is the probe restart, an empty window that scores 0. Middle: other benign hosts with a rule (exploratory, Section~\ref{sec:res-explore}). Right: share of episodes with a rule on a surge host (C1, counted before the surge as well) or on a former attacker (C2).}
\label{fig:timeseries}
\end{figure*}

\begin{table*}[!t]
\centering
\caption{Planned parameters after the event: IQM over 10 seeds [bootstrap 95\% interval]}
\label{tab:planned}
\footnotesize
\setlength{\tabcolsep}{4pt}
\begin{tabular}{lllll}
\toprule
Parameter & A2 & A1 & A3 & Static \\
\midrule
\multicolumn{5}{l}{\textit{C1: benign surge}}\\
Service after the event & 0.980 [0.980, 0.982] & 0.981 [0.980, 0.983] & 0.980 [0.980, 0.980] & 0.980 [0.980, 0.980] \\
Share of steps with a rule on a surge host & 0.563 [0.478, 0.777] & 0.560 [0.395, 0.820] & 0.000 [0.000, 0.000]$^{*}$ & 0.504 [0.481, 0.528] \\
Surge-host delivery & 0.346 [0.213, 0.405] & 0.334 [0.176, 0.479] & 0.848 [0.838, 0.861]$^{*}$ & 0.351 [0.340, 0.372] \\
Other benign delivery (post/pre) & 0.882 [0.805, 0.933] & 0.855 [0.732, 0.901] & 0.915 [0.871, 0.950] & 0.900 [0.863, 0.936] \\
Collateral & 0.629 [0.521, 0.815] & 0.627 [0.445, 0.866] & 0.000 [0.000, 0.000]$^{*}$ & 0.603 [0.586, 0.618] \\
Collateral, other benign hosts & 0.290 [0.168, 0.521] & 0.317 [0.145, 0.577] & 0.000 [0.000, 0.000]$^{*}$ & 0.313 [0.296, 0.329] \\
\midrule
\multicolumn{5}{l}{\textit{C2: attacker swap}}\\
Service after the event & 0.958 [0.950, 0.969] & 0.962 [0.954, 0.973] & 0.957 [0.928, 0.970] & 0.929 [0.917, 0.941]$^{*}$ \\
Worst window, steps 52--99 & 0.367 [0.167, 0.667] & 0.400 [0.200, 0.767] & 0.233 [0.067, 0.433] & 0.100 [0.000, 0.200] \\
Recovery latency (steps) & 0.40 [0.00, 1.40] & 1.50 [0.30, 2.40] & 2.53 [1.20, 4.47] & 2.35 [0.80, 4.03] \\
Steps with a rule on a former attacker & 16.0 [12.5, 32.9] & 19.7 [10.7, 37.8] & 6.8 [5.7, 7.3]$^{*}$ & 17.2 [15.9, 18.1] \\
Steps until no such rule & 5.87 [4.33, 27.93] & 9.93 [2.80, 29.33] & 1.23 [0.87, 2.70]$^{*}$ & 4.03 [3.40, 4.53]$^{*}$ \\
Former-attacker delivery & 0.737 [0.374, 0.820] & 0.748 [0.434, 0.891] & 0.867 [0.821, 0.914]$^{*}$ & 0.831 [0.800, 0.869] \\
Attack delivery (post/pre) & 0.760 [0.496, 0.855] & 0.727 [0.439, 0.843] & 0.844 [0.821, 0.878] & 0.741 [0.703, 0.780] \\
Other benign delivery (post/pre) & 0.862 [0.798, 0.901] & 0.838 [0.680, 0.916] & 0.891 [0.856, 0.931] & 0.884 [0.850, 0.928] \\
Collateral & 0.253 [0.193, 0.470] & 0.278 [0.164, 0.541] & 0.103 [0.088, 0.116]$^{*}$ & 0.307 [0.293, 0.324] \\
Collateral, other benign hosts & 0.136 [0.086, 0.268] & 0.146 [0.073, 0.324] & 0.043 [0.032, 0.052]$^{*}$ & 0.193 [0.179, 0.210] \\
\bottomrule
\end{tabular}
\par\smallskip\parbox{\textwidth}{\footnotesize $^{*}$Differs from A2 (exact paired Wilcoxon, Holm over the three comparisons, $p_H<0.05$). Service is probe success; collateral is the share of active rules on hosts benign at that step; delivery is a share of the traffic offered or of the level before the event. In C1 every arm had a worst window of 1.000 and a recovery latency of 0; decision cost is given in Section~\ref{sec:res-cost}.}
\end{table*}

Each value in Table~\ref{tab:planned} is an interquartile mean (IQM) over ten seeds, and $p_H$ is the Holm-adjusted $p$ of an exact paired Wilcoxon test of A2 against the named arm (Section~\ref{sec:stats}); $p_H=0.0059$ means that A2 fell on the same side of the other arm in all ten seeds. Two points of measurement matter here. The event at step 50 restarts the probe clients, and $S$ was 0 at that step in 86\% to 100\% of episodes in every arm because no probe ran and an empty window scores 0. We therefore take the worst window over steps 52 to 99 (ISASF's version, from step 2, is between 0.000 and 0.033 for every arm) and start the recovery clock at step 51, where ISASF starts it. A3 is given the true regime label, and in a benign episode that label is ``none'', on which A3 returns a no-op without running its network.

\subsection{C1: benign surge}\label{sec:res-c1}

At step 50 of a benign episode two benign hosts raise their rate to 4.0\,Mbps. No arm lost service. Mean $S$ after the event was 0.980 or 0.981 in every arm, $S$ stayed at 1 from step 52 to the end of all 200 episodes, and every episode met the recovery rule at once; the shortfall below 1 is the empty window at step 50, one step in fifty.

As Section~\ref{sec:conditions} predicted, A1, A2 and the static rule treated the surge as an attack. They kept a rule on at least one surge host in 0.563 (A2), 0.560 (A1) and 0.504 (static) of the steps after the event, and the surge hosts got 0.346, 0.334 and 0.351 of their offered traffic through (Fig.~\ref{fig:seeds}, top row). A3 never acted and let 0.848 through. A2 differs from A3 on both measures in all ten seeds ($p_H=0.0059$) and from neither A1 nor the static rule ($p_H \geq 0.50$). The other benign hosts gained nothing measurable from the blocking: their delivery after the event, relative to their own level before it, was 0.882 under A2 and 0.900 under the static rule against 0.915 under A3 ($p_H \geq 0.19$ for every comparison with A2). Collateral was 0.629 for A2, 0.627 for A1 and 0.603 for the static rule, and on the other benign hosts alone 0.290, 0.317 and 0.313, with no difference between A2 and either arm ($p_H \geq 0.86$).

\subsection{C2: attacker swap}\label{sec:res-c2}

A UDP flood from three attackers runs from step 0. At step 50 those three switch to benign traffic and three hosts that had been benign start the flood, so a rule left on a former attacker now blocks a legitimate user.

A2 kept service higher than the static rule after the swap, 0.958 against 0.929 (median paired difference 0.028). It was ahead in 8 of the 9 seeds in which the two differed ($p=0.0078$, $p_H=0.023$) and level with A1 and A3 ($p_H=1.0$). The worst window put the arms in the same order (A1 0.400, A2 0.367, A3 0.233, static 0.100) but no difference reached significance ($p_H=0.20$ against the static rule). One of the 200 episodes, under the static rule, never recovered. Recovery latency (A2 0.40 steps, A1 1.50, A3 2.53, static 2.35) did not separate A2 from any arm ($p_H=0.19$ each).

On the rules left on the former attackers, A3 does best and A2 no better than the static rule. A rule sat on at least one former attacker (Fig.~\ref{fig:timeseries}, right) in 16.0 of the 50 steps after the swap for A2, 19.7 for A1, 17.2 for the static rule and 6.8 for A3. A2 took 5.87 steps to reach a step with no such rule (an arm that never got there scores 50), the static rule 4.03 and A3 1.23. A2 is worse than A3 on both measures ($p_H=0.012$) and slower to clear than the static rule by a median of 1.9 steps ($p_H=0.043$). We had predicted the opposite in Section~\ref{sec:conditions}: the static rule clears only when no source is flagged, so we expected its rules on the former attackers to last the whole second half, but its held-over rules lasted 2.87 steps (Table~\ref{tab:explore}). The former attackers' traffic reached the victim at 0.737 of the benign level under A2 and 0.867 under A3 ($p_H=0.012$). Attack delivery after the swap relative to before (0.727 to 0.844) and delivery from the other benign hosts did not separate A2 from any arm ($p_H \geq 0.082$), and A3 kept the lowest collateral, 0.103 against 0.253 for A2 ($p_H=0.0059$).

\begin{figure*}[!t]
\centering
\includegraphics[width=\textwidth]{fig_seed_parameters}
\caption{Planned parameters per seed. Dots: one seed (mean of its 5 episodes); marker and bar: IQM over the 10 seeds with its bootstrap 95\% interval. The two A2 seeds near 50 in the bottom-right panel are training seeds 4 and 8 (Section~\ref{sec:res-explore}).}
\label{fig:seeds}
\end{figure*}

\begin{table*}[!t]
\centering
\caption{Exploratory measures, added after the planned results were seen (not pre-registered): IQM over 10 seeds}
\label{tab:explore}
\footnotesize
\setlength{\tabcolsep}{4pt}
\begin{tabular}{lllll}
\toprule
Measure & A2 & A1 & A3 & Static \\
\midrule
\multicolumn{5}{l}{\textit{C1: before (steps 2--49) $\rightarrow$ after (steps 50--99) the surge}}\\
Surge hosts, share of steps with a rule & 0.090 $\rightarrow$ 0.499$^{*}$ & 0.177 $\rightarrow$ 0.518$^{*}$ & 0.000 $\rightarrow$ 0.000 & 0.225 $\rightarrow$ 0.454$^{*}$ \\
Other benign hosts, share of steps with a rule & 0.095 $\rightarrow$ 0.155$^{*}$ & 0.193 $\rightarrow$ 0.219 & 0.000 $\rightarrow$ 0.000 & 0.219 $\rightarrow$ 0.167$^{*}$ \\
Other benign hosts with a rule, per step & 0.665 $\rightarrow$ 1.086$^{*}$ & 1.354 $\rightarrow$ 1.533 & 0.000 $\rightarrow$ 0.000 & 1.532 $\rightarrow$ 1.171$^{*}$ \\
\midrule
\multicolumn{5}{l}{\textit{C2: after the swap (steps 50--99)}}\\
Steps with a held-over rule & 4.07 & 9.07 & 0.53$^{\ddagger}$ & 2.87 \\
Steps with only re-added rules & 10.13 & 8.63 & 5.73 & 14.27$^{\ddagger}$ \\
Per-host rule rate, former attackers vs hosts benign throughout & 0.114 vs 0.116 & 0.102 vs 0.178 & 0.068 vs 0.034$^{*}$ & 0.158 vs 0.161 \\
\bottomrule
\end{tabular}
\par\smallskip\parbox{\textwidth}{\footnotesize $^{*}$Within-arm change (C1) or difference (C2: former attackers against hosts benign throughout), exact paired Wilcoxon, Holm over the arms, $p_H<0.05$. $^{\ddagger}$Differs from A2, Holm over the three comparisons, $p_H<0.05$. A held-over rule was on a former attacker at step 49 and at every step since; any other rule on a former attacker is re-added.}
\end{table*}

\subsection{Decision cost}\label{sec:res-cost}

A2 is the slowest to decide. Its median inference time per step was 0.688\,ms in C1 and 0.696\,ms in C2, against 0.323 and 0.329\,ms for A1, and it was slower than every other arm in all ten seeds of both conditions ($p_H=0.0059$), though still under 0.15\% of a 0.5\,s step. Installing the rules through the controller's RESTCONF interface took A2 a median of 3.84\,ms in C1 and 5.00\,ms in C2, less than the static rule in C1 (8.60\,ms, $p_H=0.0059$) and than A3 in C2 (15.76\,ms, $p_H=0.029$). Mean controller CPU stayed between 2.67\% and 3.84\%, and A2 loaded it less than the static rule in C1 (3.27\% against 3.58\%, $p_H=0.0078$).

\subsection{Exploratory analyses (not pre-registered)}\label{sec:res-explore}

We added the measures in Table~\ref{tab:explore} after seeing the planned results, to test two readings of them. They were not pre-registered and their $p$-values are corrected only within their own small families, so they describe these runs but confirm nothing.

\emph{C1: did the surge draw the rules?} Before the surge, the hosts that would later surge were blocked about as often as the other benign hosts: per host, 0.090 against 0.095 for A2, 0.177 against 0.193 for A1 and 0.225 against 0.219 for the static rule. Afterwards their blocking rose in all ten seeds for each of the three, to 0.499, 0.518 and 0.454 ($p_H=0.0059$ each), so the rules follow the surge. Counting the blocked bystanders changes the collateral reading of Section~\ref{sec:res-c1}. Once the surge began, A2 blocked more of the other benign hosts, from 0.665 hosts with a rule per step to 1.086 in all ten seeds ($p_H=0.0059$), while the static rule went from 1.532 to 1.171 (down in 9 of 10, $p_H=0.0078$) and A1 showed no consistent change (up in 7 of 10, $p_H=0.068$). After the event the two counts are close (1.086 against 1.171, $p_H=1.0$), as the planned share also says, but the share cannot show that A2 and the static rule got there from opposite directions.

\emph{C2: held-over or re-added rules?} We split the steps with a rule on a former attacker into those with a held-over rule and those with only re-added rules. On means, which add up where IQMs do not, A2's 21.38 steps are 12.54 held-over and 8.84 re-added, and the static rule's 16.98 are 2.88 and 14.10. A2's held-over mean is far above its IQM of 4.07 because two of its seeds (training seeds 4 and 8) kept a held-over rule for all 50 steps of every episode. On held-over rules A2 and the static rule do not differ (median paired difference 0, $p_H=0.89$), and A2 holds them longer than A3 ($p_H=0.035$). Re-added rules land on former attackers no more often than on hosts that were benign throughout (per host, 0.114 against 0.116 for A2, 0.102 against 0.178 for A1 and 0.158 against 0.161 for the static rule; $p_H=1.0$ each), so for these arms a re-added rule is an ordinary false flag. A3 is the exception, at 0.068 against 0.034 in all ten seeds ($p_H=0.0078$). Two planned C2 results read differently as a result. The stale-rule count mixes rules carried over from the flood with background false flags, mostly the latter for the static rule. A2's slower clearing than the static rule does not come from carried-over rules, on which the two are level; over the second half A2 had fewer steps with only re-added rules in all ten seeds (10.13 against 14.27, $p_H=0.0059$).
